\documentclass{article}
\usepackage{amsmath,amssymb,amsfonts,amsthm}
\usepackage{algorithm}
\usepackage{algpseudocode}
\usepackage{enumitem}
\usepackage{hyperref}
\newtheorem{theorem}{Theorem}
\newtheorem{lemma}{Lemma}
\newtheorem{assumption}{Assumption}
\newcommand{\prox}{\operatorname{prox}}
\newcommand{\grad}{\nabla}
\begin{document}

\section*{Algorithm Name}
\textbf{Accelerated Proximal Gradient Method with Momentum on the Gradient Input}  
(also known as FISTA).

\section*{Mathematical Formulation}
Let $f:\mathbb{R}^{d}\to\mathbb{R}$ be convex and $L$‑smooth (i.e. $\grad f$ is $L$‑Lipschitz) and let $g:\mathbb{R}^{d}\to\mathbb{R}\cup\{+\infty\}$ be convex, possibly non‑smooth.  
Given a stepsize $\eta\in(0,1/L]$ and a sequence $\{\beta_k\}_{k\ge 0}\subseteq[0,1)$, the iterates are defined as  

\[
\begin{aligned}
y_k &= x_k + \beta_k (x_k - x_{k-1}), \qquad (k\ge 0,\; x_{-1}=x_0)\\[4pt]
x_{k+1} &= \prox_{\eta g}\!\bigl(y_k - \eta \,\grad f(y_k)\bigr),
\end{aligned}
\]
where  

\[
\prox_{\eta g}(z) \;:=\; \arg\min_{u\in\mathbb{R}^{d}}\Bigl\{ g(u)+\frac{1}{2\eta}\|u-z\|^{2}\Bigr\}.
\]

A common choice for the momentum coefficients is the Nesterov schedule  

\[
t_0=1,\qquad 
t_{k+1}= \frac{1+\sqrt{1+4t_{k}^{2}}}{2},\qquad 
\beta_k = \frac{t_k-1}{t_{k+1}} .
\]

\section*{Pseudocode}
\begin{algorithm}[H]
\caption{Accelerated Proximal Gradient (APG)}\label{alg:apg}
\begin{algorithmic}[1]
\Require Convex $f$ with $L$‑Lipschitz gradient, convex $g$, stepsize $\eta\in(0,1/L]$, initial point $x_0\in\operatorname{dom} g$.
\State $x_{-1}\gets x_0$, $t_0\gets 1$
\For{$k=0,1,2,\dots$}
    \State $t_{k+1}\gets \dfrac{1+\sqrt{1+4t_k^2}}{2}$
    \State $\beta_k\gets \dfrac{t_k-1}{t_{k+1}}$
    \State $y_k \gets x_k + \beta_k (x_k - x_{k-1})$
    \State $x_{k+1} \gets \prox_{\eta g}\!\bigl(y_k - \eta \,\grad f(y_k)\bigr)$
    \If{termination criterion satisfied}
        \State \textbf{break}
    \EndIf
\EndFor
\State \textbf{return} $x_{k+1}$
\end{algorithmic}
\end{algorithm}

\section*{Dry Run Example}
Consider the simple composite problem  

\[
\min_{w\in\mathbb{R}} \; f(w)+g(w),\qquad 
f(w)=(w-3)^2,\;\; g(w)=0,
\]
so $\grad f(w)=2(w-3)$, $L=2$, and we choose $\eta=0.1$ (indeed $\eta\le 1/L=0.5$).  
Initialize $x_0=0$, $x_{-1}=0$, and use the Nesterov schedule.

\begin{center}
\begin{tabular}{c|c|c|c|c}
Iter $k$ & $t_k$ & $\beta_k$ & $y_k$ & $x_{k+1}$ \\ \hline
0 & $1$ & $0$ & $y_0 = 0$ & $x_1 = \prox_{0.1\cdot 0}\!\bigl(0-0.1\cdot 2(0-3)\bigr)=0+0.6=0.6$ \\
1 & $t_1=\frac{1+\sqrt{5}}{2}\approx1.618$ & $\beta_0=\frac{1-1}{1.618}=0$ & $y_1 = x_1=0.6$ &
\begin{aligned}
&\grad f(y_1)=2(0.6-3)=-4.8\\
&x_2 = 0.6-0.1(-4.8)=0.6+0.48=1.08
\end{aligned}\\
2 & $t_2=\frac{1+\sqrt{1+4(1.618)^2}}{2}\approx2.193$ & $\beta_1=\frac{1.618-1}{2.193}\approx0.282$ &
\begin{aligned}
y_2 &= x_2+0.282\,(x_2-x_1)\\
    &=1.08+0.282\,(1.08-0.6)=1.08+0.282\cdot0.48\approx1.215
\end{aligned}
&
\begin{aligned}
&\grad f(y_2)=2(1.215-3)=-3.57\\
&x_3 = 1.215-0.1(-3.57)=1.215+0.357\approx1.572
\end{aligned}
\end{tabular}
\end{center}
The objective values are $f(0.6)=5.76$, $f(1.08)=3.69$, $f(1.572)=2.06$, showing rapid decrease toward the optimum $f(3)=0$.

\section*{Assumptions}
\begin{assumption}[Smoothness of $f$]\label{ass:smooth}
$f$ is convex and differentiable with $L$‑Lipschitz gradient:
\[
\|\grad f(u)-\grad f(v)\|\le L\|u-v\|,\qquad\forall u,v\in\mathbb{R}^{d}.
\]
\end{assumption}
\begin{assumption}[Convexity of $g$]\label{ass:convex}
$g$ is proper, closed and convex (possibly non‑smooth). Its proximal operator $\prox_{\eta g}$ is well defined for any $\eta>0$.
\end{assumption}
\begin{assumption}[Stepsize]\label{ass:stepsize}
The stepsize satisfies $0<\eta\le 1/L$.
\end{assumption}
\begin{assumption}[Momentum schedule]\label{ass:beta}
The sequence $\{t_k\}$ is defined by $t_0=1$, $t_{k+1}=\frac{1+\sqrt{1+4t_k^2}}{2}$ and $\beta_k=\frac{t_k-1}{t_{k+1}}\in[0,1)$.
\end{assumption}

\section*{Main Convergence Theorem}
\begin{theorem}[Optimal $O(1/k^2)$ rate]\label{thm:rate}
Let $\{x_k\}$ be generated by Algorithm~\ref{alg:apg} under Assumptions~\ref{ass:smooth}–\ref{ass:beta}. Define the composite objective $F(x)=f(x)+g(x)$. Then for all $k\ge 1$,
\[
F(x_k)-F(x^\star)\;\le\;\frac{2\|x_0-x^\star\|^{2}}{\eta\,t_k^{2}},
\]
where $x^\star\in\arg\min F$.
Consequently, using $t_k\ge \frac{k+1}{2}$ one obtains  
\[
F(x_k)-F(x^\star)=O\!\left(\frac{1}{k^{2}}\right).
\]
\end{theorem}

\section*{Theoretical Properties}
\begin{itemize}[leftmargin=*]
\item \textbf{Stability:} The iterate sequence remains bounded because the proximal step is non‑expansive and the momentum coefficients satisfy $\beta_k<1$.
\item \textbf{Hyperparameter sensitivity:} The convergence guarantee holds for any $\eta\in(0,1/L]$. Larger $\eta$ (up to $1/L$) yields a smaller constant in the bound.
\item \textbf{Comparison to baselines:}
    \begin{itemize}
    \item Plain proximal gradient (no momentum) achieves $O(1/k)$ rate.
    \item The accelerated scheme improves the rate to $O(1/k^{2})$, which is optimal for the class of convex composite problems.
    \end{itemize}
\end{itemize}

\section*{Discussion}
The theorem shows that the accelerated proximal gradient method attains the same optimal rate as Nesterov’s accelerated gradient for smooth problems, despite the presence of a non‑smooth term $g$. The proof hinges on a potential function that couples the objective error with a weighted distance to the optimum, and on the specific choice of $\beta_k$ derived from the sequence $\{t_k\}$.

\section*{Preliminaries (Definitions and Lemmas)}
\begin{lemma}[Descent Lemma for $L$‑smooth $f$]\label{lem:descent}
For any $u,v\in\mathbb{R}^{d}$,
\[
f(v)\le f(u)+\langle\grad f(u),v-u\rangle+\frac{L}{2}\|v-u\|^{2}.
\]
\end{lemma}
\begin{lemma}[Proximal inequality]\label{lem:prox}
For any $z\in\mathbb{R}^{d}$ and any $u\in\operatorname{dom}g$,
\[
g(\prox_{\eta g}(z))+\frac{1}{2\eta}\|\prox_{\eta g}(z)-z\|^{2}
\le g(u)+\frac{1}{2\eta}\|u-z\|^{2}-\frac{1}{2\eta}\|u-\prox_{\eta g}(z)\|^{2}.
\]
\end{lemma}

\section*{Main Proof (Step‑by‑Step Derivation)}
We prove Theorem~\ref{thm:rate}.  Throughout, let $x^\star$ be a minimizer of $F$.

\noindent\textbf{Notation:} Define the auxiliary sequence $y_k$ and the potential
\[
\Phi_k \;:=\; t_k^{2}\bigl(F(x_k)-F(x^\star)\bigr) + \frac{1}{2\eta}\|x_k - x^\star\|^{2}.
\]

\noindent\underline{[Step 1]} \textbf{Apply Lemma~\ref{lem:prox} with $z = y_k-\eta\grad f(y_k)$ and $u = x^\star$.}
\[
\begin{aligned}
g(x_{k+1}) + \frac{1}{2\eta}\|x_{k+1}-y_k+\eta\grad f(y_k)\|^{2}
&\le g(x^\star)+\frac{1}{2\eta}\|x^\star-y_k+\eta\grad f(y_k)\|^{2}\\
&\quad -\frac{1}{2\eta}\|x^\star-x_{k+1}\|^{2}.
\end{aligned}
\tag{1}
\]

\noindent\underline{[Step 2]} \textbf{Add $f(x_{k+1})$ to both sides of (1).} Using the descent lemma (Lemma~\ref{lem:descent}) with $u=y_k$ and $v=x_{k+1}$,
\[
f(x_{k+1})\le f(y_k)+\langle\grad f(y_k),x_{k+1}-y_k\rangle+\frac{L}{2}\|x_{k+1}-y_k\|^{2}.
\tag{2}
\]

\noindent\underline{[Step 3]} \textbf{Combine (1) and (2).} The terms $\langle\grad f(y_k),x_{k+1}-y_k\rangle$ cancel,
\[
\begin{aligned}
F(x_{k+1}) &\le f(y_k)+g(x^\star) +\frac{1}{2\eta}\|x^\star-y_k+\eta\grad f(y_k)\|^{2} \\
&\quad -\frac{1}{2\eta}\|x^\star-x_{k+1}\|^{2}
+ \frac{L-\frac{1}{\eta}}{2}\|x_{k+1}-y_k\|^{2}.
\end{aligned}
\tag{3}
\]

\noindent\underline{[Step 4]} \textbf{Use the stepsize condition $\eta\le 1/L$.} Hence $L-\frac{1}{\eta}\le 0$ and the last term in (3) is non‑positive. Dropping it yields
\[
F(x_{k+1}) \le f(y_k)+g(x^\star) +\frac{1}{2\eta}\|x^\star-y_k+\eta\grad f(y_k)\|^{2}
-\frac{1}{2\eta}\|x^\star-x_{k+1}\|^{2}.
\tag{4}
\]

\noindent\underline{[Step 5]} \textbf{Expand the squared norm in (4).}
\[
\|x^\star-y_k+\eta\grad f(y_k)\|^{2}
= \|x^\star-y_k\|^{2} +2\eta\langle\grad f(y_k),x^\star-y_k\rangle + \eta^{2}\|\grad f(y_k)\|^{2}.
\tag{5}
\]

\noindent\underline{[Step 6]} \textbf{Apply convexity of $f$:}
\[
f(y_k)-f(x^\star)\le \langle\grad f(y_k),y_k-x^\star\rangle = -\langle\grad f(y_k),x^\star-y_k\rangle .
\tag{6}
\]

\noindent\underline{[Step 7]} \textbf{Insert (5) and (6) into (4).} After rearranging,
\[
\begin{aligned}
t_{k+1}^{2}\bigl(F(x_{k+1})-F(x^\star)\bigr)
&\le t_{k+1}^{2}\Bigl[\frac{1}{2\eta}\|x^\star-y_k\|^{2}
-\frac{1}{2\eta}\|x^\star-x_{k+1}\|^{2}\Bigr]   \\
&\quad + t_{k+1}^{2}\Bigl[\frac{\eta}{2}\|\grad f(y_k)\|^{2}\Bigr].
\end{aligned}
\tag{7}
\]

\noindent\underline{[Step 8]} \textbf{Relate $y_k$ to $x_k$ and $x_{k-1}$.} By definition,
\[
y_k = x_k + \beta_k(x_k-x_{k-1}),\qquad \beta_k=\frac{t_k-1}{t_{k+1}}.
\tag{8}
\]

\noindent\underline{[Step 9]} \textbf{Show the key recursion for the potential.}
A direct but elementary algebraic manipulation using (8) gives
\[
t_{k+1}^{2}\|x^\star-y_k\|^{2}
= t_k^{2}\|x^\star-x_k\|^{2} - (t_{k+1}^{2}-t_k^{2})\|x_k-x_{k-1}\|^{2}.
\tag{9}
\]

\noindent\underline{[Step 10]} \textbf{Insert (9) into (7) and drop the non‑negative term
$(t_{k+1}^{2}-t_k^{2})\|x_k-x_{k-1}\|^{2}$}. We obtain
\[
t_{k+1}^{2}\bigl(F(x_{k+1})-F(x^\star)\bigr)
\le t_k^{2}\bigl(F(x_k)-F(x^\star)\bigr)
+ \frac{1}{2\eta}\bigl(\|x^\star-x_k\|^{2}-\|x^\star-x_{k+1}\|^{2}\bigr).
\tag{10}
\]

\noindent\underline{[Step 11]} \textbf{Define the potential $\Phi_k$ and rewrite (10).}
Recall $\Phi_k = t_k^{2}\bigl(F(x_k)-F(x^\star)\bigr) + \frac{1}{2\eta}\|x_k-x^\star\|^{2}$.  
Equation (10) is exactly $\Phi_{k+1}\le \Phi_k$.

\noindent\underline{[Step 12]} \textbf{Monotonicity of $\Phi_k$.} Since $\Phi_{k+1}\le\Phi_k$ for all $k$, we have for any $k\ge1$,
\[
\Phi_k \le \Phi_0
= t_0^{2}\bigl(F(x_0)-F(x^\star)\bigr)+\frac{1}{2\eta}\|x_0-x^\star\|^{2}
\le \frac{1}{2\eta}\|x_0-x^\star\|^{2},
\]
because $t_0=1$ and $F(x_0)-F(x^\star)\ge0$.

\noindent\underline{[Step 13]} \textbf{Extract the bound on the function value.}
From the definition of $\Phi_k$,
\[
t_k^{2}\bigl(F(x_k)-F(x^\star)\bigr)
\le \Phi_k \le \frac{1}{2\eta}\|x_0-x^\star\|^{2},
\]
which yields the claimed inequality
\[
F(x_k)-F(x^\star)\le \frac{\|x_0-x^\star\|^{2}}{2\eta\,t_k^{2}}.
\tag{11}
\]

\noindent\underline{[Step 14]} \textbf{Lower bound on $t_k$.} By induction on the recursion $t_{k+1}= \frac{1+\sqrt{1+4t_k^{2}}}{2}$ one shows $t_k\ge\frac{k+1}{2}$. Substituting into (11) gives
\[
F(x_k)-F(x^\star)\le \frac{2\|x_0-x^\star\|^{2}}{\eta\,(k+1)^{2}} = O\!\left(\frac{1}{k^{2}}\right).
\]

Thus Theorem~\ref{thm:rate} is proved. \hfill $\square$

\section*{Rate Derivation (With Constants)}
The constant in the bound is explicitly $\dfrac{2\|x_0-x^\star\|^{2}}{\eta}$.
If the stepsize is chosen as the maximal admissible value $\eta=1/L$, the bound becomes
\[
F(x_k)-F(x^\star)\le \frac{2L\|x_0-x^\star\|^{2}}{(k+1)^{2}}.
\]

\section*{Special Cases}
\begin{enumerate}
\item \textbf{Pure smooth case ($g\equiv0$).} The algorithm reduces to Nesterov’s accelerated gradient method, and the same $O(1/k^{2})$ bound holds.
\item \textbf{Indicator regularizer $g=\iota_{\mathcal{C}}$.} The proximal step becomes Euclidean projection onto a convex set $\mathcal{C}$, yielding the accelerated projected gradient method.
\item \textbf{Strongly convex $F$.} If $F$ is $\mu$‑strongly convex, a simple modification of the potential yields a linear convergence rate $O\bigl((1-\sqrt{\mu/L})^{k}\bigr)$.
\end{enumerate}

\end{document}